% Part: first-order-logic
% Chapter: axiomatic-deduction
% Section: rules-and-proofs

\documentclass[../../../include/open-logic-section]{subfiles}

\begin{document}

\iftag{FOL}
      {\olfileid{fol}{axd}{rul}}
      {\olfileid{pl}{axd}{rul}}

\olsection{నియమాలు మరియు \usetoken{P}{derivation}}

\begin{explain}
  స్వీకృతాధారిత !!{derivation}s బహుశా తర్కంలోని అత్యంత సరళమైన
  !!{derivation} వ్యవస్థలు. ఒక !!^a{derivation} కేవలం !!{formula}s
  క్రమం. అది ఒక !!a{derivation}గా పరిగణింపబడాలంటే, ఆ క్రమంలోని ప్రతి
  !!{formula} ఒక స్వీకృత రూపమైనా కావాలి, లేదా నిగమన నియమం ద్వారా
  క్రమంలో దానికి ముందున్న ఒకటి లేదా అంతకంటే ఎక్కువ !!{formula}s నుంచి
  రావాలి. ఒక !!^a{derivation} తన చివరి !!{formula}ను !!{derive}s చేస్తుంది.
\end{explain}

\begin{defn}[!!^{derivability}]
$\Gamma$ అనేది $\Lang L$ యొక్క !!{formula}s సమితి అయితే,
$\Gamma$ నుంచి ఒక \emph{!!{derivation}} అంటే $!A_1$,
\dots,~$!A_n$ అనే !!{formula}s పరిమిత క్రమం; ప్రతి $i \le n$కు
కిందివాటిలో ఒకటి నెరవేరాలి:
\begin{enumerate}
\item $!A_i \in \Gamma$; లేదా
\item $!A_i$ ఒక స్వీకృతం; లేదా
\item $j < i$ (మరియు $k < i$) అయ్యే ఏవో $!A_j$ (మరియు $!A_k$) నుంచి
  ఒక నిగమన నియమం ద్వారా $!A_i$ వస్తుంది.
\end{enumerate}
\end{defn}

ఏది సరైన !!{derivation}గా పరిగణింపబడుతుందో, మనం ఏ నిగమన నియమాలను
అనుమతిస్తామనే దానిపై (మరియు, తప్పకుండా, వేటిని స్వీకృతాలుగా
తీసుకుంటామనే దానిపై) ఆధారపడి ఉంటుంది. ఒక నిగమన నియమం అనేది కొన్ని
షరతుల కింద !!a{derivation}లోని $A_i$ అనే దశ సరైన నిగమన దశ అని
చెప్పే “అయితే--అప్పుడు” ప్రకటన.

\begin{defn}[నిగమన నియమం]
$\Gamma$ నుంచి ఒక !!a{derivation}లో ఏది సరైన నిగమన దశగా
పరిగణింపబడుతుందో దానికి ఒక \emph{నిగమన నియమం} సరిపోయే షరతును ఇస్తుంది.
\end{defn}

ఉదాహరణకు, $!A \in \Gamma$ అయిన ఏక-మూలక క్రమం $!A$ అలవోకగా ఒక
!!a{derivation} అవుతుంది కనుక, కిందిది చాలా సరళమైన నిగమన నియమం
కావచ్చు:
\begin{quote}
  $!A \in \Gamma$ అయితే, $\Gamma$ నుంచి ఏ !!{derivation}లోనైనా
  $!A$ ఎల్లప్పుడూ ఒక సరైన నిగమన దశ.
\end{quote}
అలాగే, $!A$ స్వీకృతాల్లో ఒకటైతే, $!A$ తనంతట తానే ఒక
!!a{derivation}; కాబట్టి ఇది కూడా ఒక నిగమన నియమం:
\begin{quote}
  $!A$ ఒక స్వీకృతమైతే, $!A$ ఒక సరైన నిగమన దశ.
\end{quote}
పరిశీలిస్తున్న దశకు ముందు కనిపించే !!{formula}sను నిగమన నియమం
ఆశ్రయించినప్పుడు విషయం మరింత ఆసక్తికరంగా ఉంటుంది. కింది నియమాన్ని
\emph{మోడస్ పోనెన్స్} అంటారు:
\begin{quote}
  !!{derivation}లో ముందుగా $!B \lif !A$, $!B$ కనిపిస్తే,
  $!A$ ఒక సరైన నిగమన దశ.
\end{quote}
ఇదొక్కటే నిగమన నియమమైతే, పైన ఇచ్చిన మన !!{derivation} నిర్వచనం
ఇలా అవుతుంది: ప్రతి $i \le n$కు కిందివాటిలో ఒకటి నెరవేరితే, అప్పుడు
మాత్రమే $!A_1$, \dots,~$!A_n$ ఒక !!a{derivation}:
\begin{enumerate}
\item $!A_i \in \Gamma$; లేదా
\item $!A_i$ ఒక స్వీకృతం; లేదా
\item ఏదో $j < i$కు $!A_j$ అనేది $!B \lif !A_i$, అలాగే ఏదో $k < i$కు
  $!A_k$ అనేది~$!B$.
\end{enumerate}
చివరి ఉపవాక్యం, $!A_j$ ($!B \lif !A_i$), $!A_k$ ($!B$)ల నుంచి
మోడస్ పోనెన్స్ ద్వారా $!A_i$ వస్తుందని చెబుతుంది. మనం $1$ నుంచి~$n$
వరకు వెళ్లగలిగి, ప్రతిసారీ $\Gamma$లో ఉన్నదైనా, ఒక స్వీకృతమైనా,
లేదా ఒక నిగమన నియమం సరైన నిగమన దశ అని చెప్పేదైనా అయిన
!!a{formula}~$!A_i$ను కనుగొంటే, మొత్తం క్రమం ఒక సరైన
!!{derivation}గా పరిగణింపబడుతుంది.

\begin{defn}[!!^{derivability}]
$\Gamma$ నుంచి $!A$తో ముగిసే ఒక !!a{derivation} ఉంటే,
!!{formula}~$!A$ను $\Gamma$ నుంచి \emph{!!{derivable}} అంటాం;
దీన్ని $\Gamma \Proves !A$ అని రాస్తాం.
\end{defn}

\begin{defn}[సిద్ధాంతాలు]
ఖాళీ సమితి నుంచి $!A$కు ఒక !!a{derivation} ఉంటే, !!^a{formula}~$!A$
ఒక \emph{సిద్ధాంతం}. $!A$ సిద్ధాంతమైతే $\Proves !A$ అని,
కాకపోతే $\Proves/ !A$ అని రాస్తాం.
\end{defn}


\end{document}
